% !TEX root = ../tfg_etsiinf_NombreApellidos.tex
\chapter{Diseño de la arquitectura} \label{chp:diseno}

Este capítulo presenta el diseño de la arquitectura desarrollada para integrar
una interfaz de aprendizaje por refuerzo orientada al control de un UAV en
Aerostack2. La descripción se centra en la arquitectura y en las decisiones que permiten integrar Gymnasium con la
interfaz robótica de ROS 2. No se aborda aquí la programación detallada de los
módulos, sino las decisiones de diseño que estructuran la solución y delimitan su
comportamiento.

\section{Arquitectura general} \label{sct:diseno:arquitectura}

La arquitectura se organiza como una cadena de responsabilidades que separa el
entrenamiento del agente, la interfaz Gymnasium y la comunicación con el
simulador. En el nivel superior se sitúan los scripts de entrenamiento, encargados
de instanciar el entorno y conectar el algoritmo de aprendizaje por refuerzo con
la API estándar de Gymnasium. Esta API actúa como frontera entre el agente y el sistema robótico, de modo que
el agente no opera directamente sobre ROS 2, sino sobre un contrato de
observaciones, acciones, recompensas y condiciones de finalización.

El componente central de esta frontera es la clase \texttt{AS2TestEnv}, definida
como una especialización de \texttt{gym.Env} y registrada como
\texttt{AS2TestEnv-v0}. Su función de diseño consiste en traducir cada paso del
contrato Gymnasium en órdenes comprensibles por Aerostack2 y, a la inversa, convertir el
estado relativo del UAV en una observación apta para el agente. La comunicación
con Aerostack2 se realiza mediante \texttt{as2\_python\_api.DroneInterface},
mientras que el envío de velocidades se encapsula a través de
\texttt{SpeedMotion} y del método
\texttt{send\_speed\_command\_with\_yaw\_speed}. De este modo, el agente queda
desacoplado de los detalles de bajo nivel de ROS 2 y del simulador multirrotor de
Aerostack2 \cite{fernandez2023aerostack2}.

\begin{figure}[t]
\centering
\shorthandoff{<>}
\begin{tikzpicture}[
    node distance=1.25cm,
    block/.style={
        rectangle,
        rounded corners,
        draw,
        align=center,
        minimum width=3.2cm,
        minimum height=0.9cm,
        font=\small
    },
    arrow/.style={-{Latex[length=2mm]}, thick}
]

\node[block] (train) {Scripts de\\entrenamiento o prueba};
\node[block, below=of train] (gym) {API Gymnasium\\\texttt{reset}/\texttt{step}/\texttt{close}};
\node[block, below=of gym] (env) {\texttt{AS2TestEnv}\\Entorno de adaptación};
\node[block, below=of env] (as2api) {\texttt{DroneInterface}\\\texttt{SpeedMotion}};
\node[block, below=of as2api] (ros) {ROS 2\\comunicación};
\node[block, below=of ros] (sim) {AS2 Multirotor\\Simulator\\y overlay de reset};

\draw[arrow] (train) -- (gym);
\draw[arrow] (gym) -- (env);
\draw[arrow] (env) -- (as2api);
\draw[arrow] (as2api) -- (ros);
\draw[arrow] (ros) -- (sim);

\node[draw, dashed, rounded corners, fit=(gym)(env), inner sep=0.25cm,
      label={[font=\small]right:Frontera RL--robótica}] {};

\end{tikzpicture}
\shorthandon{<>}
\caption{Arquitectura general desarrollada, separando la lógica de aprendizaje,
la interfaz Gymnasium, la comunicación robótica con Aerostack2 y el simulador
con servicio de reinicio.}
\label{fig:diseno:arquitectura_general}
\end{figure}

La Figura~\ref{fig:diseno:arquitectura_general} resume esta separación por
capas. La interacción avanza desde los scripts de entrenamiento hacia la API
Gymnasium, la clase \texttt{AS2TestEnv}, la interfaz
\texttt{DroneInterface}/\texttt{SpeedMotion}, ROS 2 y, finalmente, el AS2
Multirotor Simulator. La última capa incorpora, en la versión final de la arquitectura,
un overlay con un servicio explícito de reinicio de estado. Esta estructura sigue
un criterio de aislamiento. El entrenamiento decide acciones en un espacio
matemático acotado, la interfaz Gymnasium aplica las transformaciones necesarias y
Aerostack2 ejecuta los comandos en el simulador. Además, la inicialización de ROS
2 se diseña como un recurso único por proceso mediante las variables de control
\texttt{\_rclpy\_initialized} y \texttt{\_rclpy\_lock}, con el fin de evitar
inicializaciones repetidas y reducir riesgos de carrera cuando se crean varias
instancias del entorno.

La incorporación del servicio de reinicio supone una parte crítica del
diseño. El reinicio deja de depender únicamente de comandos de velocidad o de
comportamientos de navegación de Aerostack2, y pasa a contar con una operación
propia del simulador para restaurar pose, velocidades y referencias internas. 
Esta decisión se introduce porque, en aprendizaje
por refuerzo, un episodio solo es válido si el estado inicial y la accionabilidad
posterior al reset son verificables. El entorno decide cuándo reiniciar, el
simulador aplica la mutación física del estado y el propio entorno vuelve a
certificar telemetría y movimiento antes de entregar la siguiente transición a
PPO.

\section{Diseño de la interfaz Gymnasium} \label{sct:diseno:gymnasium}

La interfaz Gymnasium se diseña como una adaptación entre el ciclo estándar
\texttt{reset}-\texttt{step}-\texttt{close} y las operaciones necesarias para
preparar, gobernar y finalizar el vuelo simulado. Esta decisión permite que el
agente interactúe con el UAV mediante una interfaz conocida en aprendizaje por
refuerzo, mientras que la lógica específica de Aerostack2 permanece encapsulada
dentro de este componente de la arquitectura.

El método \texttt{reset} establece el estado inicial de un episodio. Su diseño
incluye el armado del UAV, la activación del modo \emph{offboard}, el despegue y
la devolución de la primera observación junto con la información auxiliar del
episodio. Por tanto, el reinicio no se limita a reinicializar variables internas,
sino que prepara al sistema simulado para aceptar comandos de velocidad desde el
agente. Esta preparación es relevante porque el contrato de Gymnasium presupone
que, tras \texttt{reset}, la interfaz está lista para ejecutar \texttt{step} sin
operaciones adicionales por parte del algoritmo.

El método \texttt{step} constituye el núcleo de la interacción. En cada llamada,
el entorno recibe la acción continua del agente, la transforma en una orden de
velocidad enviada mediante \texttt{SpeedMotion}, actualiza la observación relativa
y calcula la recompensa y los indicadores de fin de episodio. El resultado sigue el contrato de Gymnasium y devuelve observación, recompensa,
indicador de terminación, indicador de truncamiento e información auxiliar. La
separación entre \emph{terminación} y \emph{truncamiento} permite distinguir un episodio que acaba
por alcanzar una condición del problema de otro que finaliza por límite externo de
pasos.

El método \texttt{close} se diseña como una operación segura de cierre del
entorno. Su responsabilidad es aterrizar el UAV y liberar o cerrar los recursos
asociados a la comunicación con Aerostack2. Esta decisión busca reducir el
riesgo de que el final de un experimento deje el simulador en un estado
inconsistente y concentra la lógica de apagado en el mismo componente que
gestiona la interacción durante el episodio.

El entorno contempla también un modo de diseño denominado
\texttt{randomize\_hover\_start}. Cuando está activo, el inicio, el objetivo y la
guiñada del episodio se muestrean dentro de límites configurables y se evitan
casos triviales mediante una separación mínima entre inicio y objetivo. Además se incorpora
 un margen de muestreo respecto a
las fronteras de la escena y un límite máximo de distancia inicio--objetivo. Esta
decisión evita que las observaciones relativas se saturen en $[-1,1]$ desde el
inicio del episodio y reduce la probabilidad de que un agente no entrenado
genere inmediatamente una terminación fuera de límites.

\begin{figure}[t]
\centering
\shorthandoff{<>}
\begin{tikzpicture}[
    node distance=1.4cm and 1.7cm,
    block/.style={
        rectangle,
        rounded corners,
        draw,
        align=center,
        minimum width=2.8cm,
        minimum height=0.85cm,
        font=\small
    },
    arrow/.style={-{Latex[length=2mm]}, thick}
]

\node[block] (agent) {Agente RL};
\node[block, right=of agent] (action) {Acción\\$[v_x,v_y,v_z,v_{\mathrm{yaw}}]$};
\node[block, right=of action] (env) {\texttt{AS2TestEnv}};
\node[block, below=of env] (sim) {Aerostack2\\Simulador};
\node[block, left=of sim] (obs) {Observación\\$[d_x,d_y,d_z,d_{\mathrm{yaw}}]$};
\node[block, left=of obs] (reward) {Recompensa\\y fin de episodio};

\draw[arrow] (agent) -- (action);
\draw[arrow] (action) -- (env);
\draw[arrow] (env) -- node[right, font=\scriptsize]{comando de velocidad} (sim);
\draw[arrow] (sim) -- node[below, align=center, font=\scriptsize]{estado\\del UAV} (obs);
\draw[arrow] (obs) -- (reward);
\draw[arrow] (reward) -- (agent);

\node[block, above=of env] (reset) {\texttt{reset()}\\armado, offboard, despegue};
\node[block, above=of action] (close) {\texttt{close()}\\aterrizaje y cierre};

\draw[arrow, dashed] (reset) -- (env);
\draw[arrow, dashed] (env) -- (close);

\end{tikzpicture}
\shorthandon{<>}
\caption{Ciclo de interacción entre el agente, la interfaz Gymnasium y el
simulador durante un episodio.}
\label{fig:diseno:ciclo_gymnasium}
\end{figure}

La Figura~\ref{fig:diseno:ciclo_gymnasium} muestra el ciclo de interacción que
se repite durante un episodio. A partir de este ciclo, el
Cuadro~\ref{tab:diseno:contrato_entorno} sintetiza el contrato principal que la
interfaz ofrece al algoritmo de aprendizaje.

\begin{table}[h]
    \centering
    \begin{tabular}{ll}
        \hline
        Elemento & Diseño adoptado \\
        \hline
        Observación & Vector relativo normalizado $[d_x,d_y,d_z,d_{\mathrm{yaw}}]$ \\
        Acción & Velocidades continuas $[v_x,v_y,v_z,v_{\mathrm{yaw}}]$ \\
        Recompensa & Distancia normalizada al objetivo y bonificación terminal \\
        Finalización & Éxito, salida de límites, baja altitud insegura o truncamiento temporal \\
        \hline
    \end{tabular}
    \caption{Resumen del contrato principal de la interfaz Gymnasium.}
    \label{tab:diseno:contrato_entorno}
\end{table}



\section{Diseño del reinicio de episodios y accionabilidad} \label{sct:diseno:reset_accionabilidad}

El reinicio de episodios constituye una decisión de diseño central en la arquitectura
de aprendizaje por refuerzo para robótica aérea. A diferencia de una interfaz
puramente numérico, no basta con asignar un nuevo estado inicial se debe ser
coherentes el estado físico del multirrotor, el controlador, la plataforma y las
lecturas que recibe la interfaz. Por este motivo, \texttt{reset()} se define como
una precondición de la interacción posterior y no como una operación auxiliar.

El diseño combina una pose inicial fija, o \texttt{fixed\_start\_pose}, con dos
vías de recuperación. La vía ordinaria devuelve al UAV que ya está en vuelo
mediante comandos de velocidad acotados, evitando introducir operaciones
bloqueantes entre episodios. Cuando esa maniobra no puede garantizar un estado
válido, se recurre a un reinicio gestionado por el simulador. Esta segunda vía no
se concibe como una recolocación ciega, sino como una recuperación capaz de
restaurar de forma coherente el estado necesario para continuar el control.

El criterio de aceptación es independiente de la vía empleada, el UAV debe
alcanzar y mantener la pose inicial, y el canal de control debe producir un
movimiento físico medible. La aceptación de una orden por la interfaz robótica
no basta para demostrar su efectividad. Si la interfaz generase observaciones y
recompensas mientras las acciones no modifican el estado del vehículo, la
experiencia dejaría de representar la tarea de aprendizaje. Por ello, los fallos
de reinicio o de accionabilidad se tratan como condiciones de rechazo temprano,
con diagnóstico explícito, y no como episodios válidos. La materialización de
esta estrategia en el simulador y en \texttt{AS2TestEnv} se describe en la
Sección~\ref{sct:implementacion:diagnostico_reset}.

\begin{table}[H]
\centering
\small
\setlength{\tabcolsep}{4pt}
\renewcommand{\arraystretch}{1.15}

\begin{tabular}{
    >{\raggedright\arraybackslash}p{0.25\textwidth}
    >{\raggedright\arraybackslash}p{0.35\textwidth}
    >{\raggedright\arraybackslash}p{0.30\textwidth}
}
\hline
\textbf{Estado o fallo} &
\textbf{Significado de diseño} &
\textbf{Respuesta prevista} \\
\hline

\path{reset_timeout} &
El UAV no alcanza la pose inicial dentro del tiempo definido. &
Fallo acotado del reinicio y rechazo de continuar el entrenamiento. \\

\path{post_controller_reacquire_timeout} &
Tras recrear o estabilizar el controlador, no se recupera la tolerancia de posición. &
Diagnóstico específico del acoplamiento entre reset y controlador. \\

\path{post_controller_hold_timeout} &
La pose se alcanza de forma transitoria, pero no se mantiene estable. &
Rechazo del reset para evitar datos de episodio físicamente inválidos. \\

Comando aceptado sin movimiento &
La interfaz acepta la orden, pero no se observa desplazamiento físico suficiente. &
Fallo cerrado de accionabilidad antes de iniciar o continuar el entrenamiento. \\

Servicio de reset aplicado sin accionabilidad &
El simulador informa que la pose se ha restaurado, pero el camino de comandos posterior no produce desplazamiento suficiente. &
Reintento completo del servicio y, si persiste, parada explícita del entrenamiento. \\

\hline
\end{tabular}

\caption{Estados y fallos considerados en el diseño del reinicio de episodios.}
\label{tab:diseno:fallos_reset}
\end{table}

El Cuadro~\ref{tab:diseno:fallos_reset} resume los principales estados de fallo
que guían el diseño. En todos los casos se favorece un fallo explícito y
trazable frente a un bloqueo indefinido o a la continuación de un entrenamiento
con datos contaminados por un estado físico no válido.


\section{Diseño del protocolo de evaluación de políticas} \label{sct:diseno:evaluacion}

La evaluación de una política aprendida exige que los episodios se ejecuten bajo
condiciones reproducibles. El protocolo se diseña en torno a dos decisiones. La
primera es la reproducibilidad por semillas: el episodio $i$ se genera
reiniciando el entorno con la semilla $s_0 + i$, de modo que la secuencia de
pares inicio--objetivo puede repetirse en evaluaciones posteriores. La segunda
es fijar una interfaz de evaluación común para una política PPO en sus modos
determinista y estocástico, que transforma la misma observación relativa
normalizada en acciones dentro del espacio acotado del entorno. Cada evaluación
registra métricas por episodio, como éxito, razón de terminación, pasos,
distancia final, longitud de trayectoria, error de guiñada terminal y normas de
acción, y las agrega globalmente y por tramos de distancia inicial. Así se
caracteriza el comportamiento de la política bajo condiciones cuantitativas
homogéneas.
